\documentclass[10pt,a4paper]{amsart}
\usepackage[margin=19mm]{geometry}
\usepackage{amsmath,amssymb,mathtools}
\usepackage[scr=rsfs,cal=euler]{mathalfa}
\usepackage{xfrac}
\usepackage[unicode,hidelinks,bookmarks=false]{hyperref}
\newcommand{\ie}{i.e.\ }
\newcommand{\cf}{cf.\ }
\newcommand{\resp}{respectively\ }
\newcommand{\Subsection}{\subsection}
\providecommand{\nobreakdash}{\nobreak\hbox{-}}
\setlength{\emergencystretch}{2em}
\multlinegap0pt
\usepackage{amssymb}
\usepackage[shortlabels]{enumitem}
\usepackage{braket}
\usepackage{mathtools}
\usepackage{tikz}
\newtheorem{theorem}{Theorem}
\title{A generalization of $q$-deformation of graphic arrangements to simplicial complexes}
\author{Tongyu Nian}
\date{Source: arXiv:2601.06546v2 (CC BY 4.0)}
\begin{document}
\maketitle

\noindent\emph{Source and licence.} A generalization of $q$-deformation of graphic arrangements to simplicial complexes. Version arXiv:2601.06546v2 (CC BY 4.0), distributed by arXiv under the Creative Commons Attribution 4.0 International licence, http://creativecommons.org/licenses/by/4.0/, as recorded in the arXiv licence metadata for this exact version and shown on the abstract page https://arxiv.org/abs/2601.06546v2. Full licence notice: \texttt{licenses/arxiv-2601.06546-CC-BY-4.0.txt}.\par\smallskip

\noindent\textit{Editorial scope.} Counted unit: the article's theorem filtration (source0.tex lines 406--410). The article's complete original proof of it is delivered with it (source0.tex lines 412--414, one \texttt{proof} environment of 3 lines). No mathematics is written, repaired or re-derived: the statement and the proof are the article's own lines, in the article's own notation, and no retained line is substituted or re-flowed.  No further source-local context is needed: every cross-reference of every retained line resolves inside the two delivered spans.  Nothing was omitted from any retained span, so the package carries no omission record.  No retained line contains a citation, so the wrapper carries no bibliography and no bibliography entry.  No retained line contains a figure environment, an image-inclusion command or drawing code, so no image is delivered and the package declares no asset dependency.  Editorial formatting only, in addition: an AMS \texttt{amsart} preamble that restates the article's own declarations and macros used by the retained lines, namely the environments \texttt{theorem} on separate counters, together with the standard packages those lines need.  The retained environments print in this document as Theorem 1 in order of appearance, whereas the article prints the same items with the numbering of its own full document.  The complete native source of the article as distributed in the arXiv e-print for this exact version is bundled unchanged as \texttt{sources/192420/source0.tex}.
\begin{theorem}\label{filtration}
    Let $G$ be a chordal graph with perfect elimination ordering $(v_1,v_2,\cdots,v_\ell)$, then the arrangement $\mathcal{M}(G,r)$ is supersolvable with filtration
    $$\mathcal{A}_i=\mathcal{M}({G_i},r)$$
    where $G_i$ is the subgraph induced by $\{v_1,v_2,\cdots,v_i\}$.
\end{theorem}

\begin{proof}
    Directly check the definition. Choose a pair of hyperplanes in $\mathcal{A}_r\setminus\mathcal{A}_{r-1}$, say $x_r-z^{r_i}x_i=0$ and $x_r-z^{r_j}x_j=0$, where $v_i$ and $v_j$ are vertices adjacent to $v_r$ such that $i,j<r$. By chordality of the graph, $(v_i,v_j)\in E(G)$, hence $x_i-z^{r_i-r_j}x_j=0$ is a hyperplane in $\mathcal{A}_{r-1}$ containing the intersection of both hyperplanes, satisfying the definition of filtration of a supersolvable arrangement.
\end{proof}

\end{document}
